\documentstyle[%


righttag%


,11pt%


,amssymb%


,russian%


,izvru%


]{amsart}





\newcommand{\tl}[3]{\medskip\noindent{\bf#1}\ #2 \dotfill #3 \par} 


\pagestyle{empty}


% \pagestyle{plain}


\begin{document}


% \setcounter{page}{80}


\par


\par


\centerline{\Large СОДЕРЖАНИЕ}


\bigskip


\bigskip


%%%%%%%%%% Use \newline %%%%%%%%%%%





\tl{Аксентьева Е.П.}{Степенная задача Римана с постоянными


показателями для двоякопери-\newline одических функций с нулями на контуре}{3}


\tl{Александров А.Ю.}{К вопросу об устойчивости


по неавтономному первому приближению}{13}


\tl{Кац Б.А.}{Интеграл Стильтьеса по фрактальному контуру


и некоторые его приложения}{21}


\tl{Лукашенко Т.П.}{О свойствах обобщенных систем


разложения, подобных ортогональным}{33}


\tl{Насибов Ф.Г.}{Об экстремальных задачах в классах целых  


функций, ограниченных на\newline вещественной оси}{49}


\tl{Плещинский Н.Б., Гусенкова А.А.}{Комплексные потенциалы


с логарифмическими\newline особенностями в ядрах


для упругих тел с дефектом вдоль гладкой дуги}{57}


\tl{Ушаков В.И., Иванова М.В.}


{Свойства функции Грина третьей смешанной задачи для\newline


параболического уравнения в нецилиндрической области}{68}


\bigskip





\bigskip


\par


\centerline{КРАТКИЕ СООБЩЕНИЯ}


\bigskip





\tl{Какичев В.А., Нгуен Суан Тхао.}


{Обобщенные свертки $H$-преобразований}{79}








\vfill


\noindent\raisebox{2pt}{\copyright} Казанский госудаpственный унивеpситет


\end{document}


